\documentclass[10pt]{article}
\usepackage[margin=1in]{geometry}
\usepackage[T1]{fontenc}
\usepackage{lmodern,microtype,amsmath,amssymb,amsthm}
\theoremstyle{definition}
\newtheorem{assumption}{Assumption}
\begin{document}
\section*{Reverse Fisher Gradient Descent}
We study gradient descent on the reverse Fisher loss $F$ between an equally weighted Gaussian mixture with $m\geq1$ means in $\mathbb{R}^d$, $d\geq1$, and a Gaussian target with mean $x^\star\in\mathbb{R}^d$. Both share an unknown covariance $\Sigma$ satisfying $\lambda_-I\preceq\Sigma\preceq\lambda_+I$, where $0<\lambda_-\leq\lambda_+$. The loss averages the squared Euclidean difference of their log-density gradients under the mixture. Each given initial mean $\theta_i^0$ lies within distance $R\geq0$ of $x^\star$. The oracle returns $\nabla F(\theta)$ at $\theta$. For $\varepsilon>0$, the final means must satisfy $\|\bar\theta_i-x^\star\|^2\leq\varepsilon$ for every component and every realization of algorithm randomness.

Write $E=(\mathbb{R}^d)^m$, equipped with the product inner product and norm
\[
\|\theta\|_E^2=\sum_{i=1}^m\|\theta_i\|^2.
\]
The model consists of an invertible linear map $A:\mathbb{R}^d\to\mathbb{R}^d$ and the target mean $x^\star$, with $\Sigma=AA^\top$. Define the unnormalized log densities
\[
\begin{aligned}
h_\theta(x)&=\log\left(\sum_{i=1}^m
 \exp\left(-\frac{\|A^{-1}(x-\theta_i)\|^2}{2}\right)\right),\\
h_\star(x)&=-\frac{\|A^{-1}(x-x^\star)\|^2}{2}.
\end{aligned}
\]
Their gradients with respect to $x$ equal the corresponding log-density gradients. For $Z\sim\mathcal{N}(0,I_d)$, define
\[
F(\theta)=\frac{1}{m}\sum_{\ell=1}^m
\mathbb{E}\left[
\left\|\nabla h_\theta(\theta_\ell+AZ)
-\nabla h_\star(\theta_\ell+AZ)\right\|^2
\right].
\]
The optimization problem is
\[
\min_{\theta\in E}F(\theta).
\]
The algorithm receives $d,m,\lambda_-,\lambda_+,R,\varepsilon$, and $\theta^0\in E$; $A$ and $x^\star$ are accessible only through the oracle.

Our assumptions:
\begin{assumption}[Covariance bounds]
For every $v\in\mathbb{R}^d$,
\[
\lambda_-\|v\|^2\leq\|A^\top v\|^2
\leq\lambda_+\|v\|^2.
\]
\end{assumption}

\begin{assumption}[Initial localization]
The given initial means satisfy
\[
\|\theta_i^0-x^\star\|\leq R,
\qquad i=1,\ldots,m.
\]
\end{assumption}

Our formulation is algorithm-specific: we analyze gradient descent with positive, potentially adaptive step sizes. Queries may be made at any point of $E$, and $\nabla F$ denotes the gradient under its product inner product. Starting from the given $\theta^0$, each oracle call and update have the form
\[
\begin{aligned}
g^t&=\nabla F(\theta^t),\\
\theta^{t+1}&=\theta^t-\eta_t g^t,
\qquad \eta_t>0.
\end{aligned}
\]
The step size is selected from the algorithm's current information before receiving $g^t$ and may depend on its randomness. The expectation defining $F$ is part of the objective; oracle replies are exact gradients of that objective.

After $K$ oracle calls, the output is the last iterate $\bar\theta=\theta^K$. The required accuracy is
\[
\|\bar\theta_i-x^\star\|^2\leq\varepsilon,
\qquad i=1,\ldots,m,
\]
for every realization of the algorithm's randomness.
\end{document}
